\section{Further small graphs}\label{sec:small}

\subsection{Chromatic number four at a prescribed compatibility parameter}\label{sec:stretch}

We recall the graph $F_{22}$ and its coloring from~\cite[Proposition~6.1]{Z}. It has a \emph{center} $o$ adjacent to nine \emph{terminals} $A,\bar A,B,\bar B,C,\bar C,X,Y,Z$, the three edges $A\bar A$, $B\bar B$, $C\bar C$, and six \emph{private edges} $u_rv_r$, where $u_r$ is adjacent to the two terminals of the left pair and $v_r$ to the two terminals of the right pair in row $r$ of the table
\[
\begin{array}{c|cc}
 r&\text{left}&\text{right}\\ \hline
 1&A,X&Z,\bar C\\
 2&\bar A,X&C,Z\\
 3&X,B&Y,\bar C
\end{array}
\qquad\qquad
\begin{array}{c|cc}
 r&\text{left}&\text{right}\\ \hline
 4&X,\bar B&Y,C\\
 5&A,Y&B,Z\\
 6&\bar A,Y&\bar B,Z
\end{array}
\]
With the Pauli matrices $x,y,z$, the unbarred terminals carry the two-qubit observables
\[
 A=I\otimes z,\quad B=x\otimes x,\quad C=z\otimes x,\quad X=x\otimes I,\quad Y=z\otimes I,\quad Z=y\otimes y,
\]
and a barred terminal $\bar S$ carries $-S$. The two observables of every pair in the table commute, and $AX=CZ$, $XB=YC$, $AY=-BZ$. The quantum $3$-coloring with colors $0,1,2$ is as follows. The center returns $0$. A terminal with observable $S$ returns $1$ on the eigenspace of $S$ for $+1$ and $2$ on the eigenspace for $-1$. An end of a private edge whose pair carries the observables $U,V$ returns $0,1,2$ on the ranges of
\begin{equation}\label{eq:private}
 \tfrac12(I-UV),\qquad\tfrac14(I-U)(I-V),\qquad\tfrac14(I+U)(I+V).
\end{equation}
By the three relations, the products $UV$ at the two ends of a private edge are negatives of each other.

\begin{proposition}[Stretching $F_{22}$]\label{prop:stretch}
Let $R\ge2$ and $l=\lceil(R-1)/2\rceil$. There is a graph $F_{22}^{(l)}$ with $22+69\,l$ vertices and $42+115\,l$ edges, chromatic number four and quantum chromatic number three, that has an $R$-compatible quantum $3$-coloring in local dimension four. Every quantum $3$-coloring of it has local dimension divisible by four. In particular $F_{22}^{(1)}$, with $91$ vertices and $157$ edges, serves $R=2$ and $R=3$.
\end{proposition}
\begin{proof}
\emph{Contexts.} We call a set of vertices of $F_{22}$ whose measurements commute pairwise a \emph{context}, the usual word for a set of jointly measurable observables. Up to bars, the pairs of the table are $AX$, $XB$, $BZ$, $ZC$, $CY$, $YA$, each occurring in two rows. We use twelve contexts. For each of the six pairs take its two terminals, the barred partner of its terminal in $\{A,B,C\}$, and the two ends of private edges whose pairs consist of these terminals; to the contexts of $AX$, $BZ$ and $CY$ add the center. The other six contexts are the pairs $\{u_r,v_r\}$. The measurements of one of the first six contexts are functions of two commuting observables, and the measurement of $o$ is trivial, so they commute. For a private edge with product $M=UV$ at $u_r$, the projections of $u_r$ are $\frac12(I-M)$ and two projections below $\frac12(I+M)$ by~\eqref{eq:private}, and those of $v_r$ are $\frac12(I+M)$ and two projections below $\frac12(I-M)$; any projection of the first kind and any of the second are orthogonal or comparable, hence commute. Every edge of $F_{22}$ lies in a context: the contexts of $AX$, $BZ$, $CY$ contain $o$ and together all nine terminals, the edge $A\bar A$ lies in the context of $AX$, and likewise $B\bar B$ and $C\bar C$; an edge from an end of a private edge to a terminal lies in the context of its pair; and a private edge is a context. The center lies in three contexts and every other vertex in two.

\emph{The graph.} For every context $\mathfrak c$ and every vertex $x\in\mathfrak c$ take a vertex $(x,\mathfrak c)$, a \emph{copy} of $x$; there are $3+2\cdot21=45$ copies. For every edge $xy$ of $F_{22}$ choose one context that contains it and join the copies of $x$ and $y$ in that context. A \emph{diamond} is the graph $K_4$ with one edge removed, and the two ends of the removed edge are its \emph{tips}. A \emph{chain} of $l$ diamonds consists of $l$ diamonds in a row, consecutive ones sharing a tip; its two \emph{ends} are the tips at its two extremities. It has $3l+1$ vertices and $5l$ edges, and its ends are at distance $2l$. For every vertex $x$ of $F_{22}$ join the copies of $x$ by chains of $l$ diamonds, two chains in a row for the three copies of $o$ and one chain for the two copies of any other vertex, the copies being the ends of the chains. This uses $23$ chains and gives a graph $F_{22}^{(l)}$ with $45+23(3l-1)=22+69\,l$ vertices and $42+23\cdot5l$ edges.

\emph{The coloring.} Every copy of $x$ and every tip on a chain between copies of $x$ carries the measurement $(P_0,P_1,P_2)$ of $x$, and the two other vertices of each diamond carry $(P_1,P_2,P_0)$ and $(P_2,P_0,P_1)$. The edge rule holds between copies as in $F_{22}$, and on the edges of a diamond because projectors for different outcomes of one measurement are orthogonal. The local dimension is four. A diamond contains a triangle, so the quantum chromatic number is three by~(Q3).

\emph{Chromatic number.} In a proper $3$-coloring of a diamond the two tips have the same color. A proper $3$-coloring of $F_{22}^{(l)}$ would thus give all copies of each vertex $x$ of $F_{22}$ the same color, and these colors would form a proper $3$-coloring of $F_{22}$, which does not exist. A proper $4$-coloring of $F_{22}$ extends: copies and tips receive the color of their vertex, and the two other vertices of a diamond two of the remaining colors.

\emph{Compatibility.} Since $2l+1\ge R$, it suffices to show that the measurements at the two ends of a path $\gamma$ with at most $2l+1$ edges commute. All vertices of a chain between copies of $x$ carry the measurement of $x$ up to a permutation of the outcomes, and every edge of the graph lies in a chain or joins two copies in one context. If the copies on $\gamma$ lie in at most one context $\mathfrak c$, then every vertex of $\gamma$ is a copy in $\mathfrak c$ or lies on a chain with an end in $\mathfrak c$, or $\gamma$ lies in a single chain; in each case the measurements along $\gamma$ are, up to permutations of outcomes, measurements of one context, or of one vertex, and commute. Otherwise $\gamma$ contains a subpath that leads from a copy in one context to a copy in another one with no copy in between. This subpath runs through a whole chain, between copies of some vertex $x$, and has at least $2l$ edges. So $\gamma$ has at most one further edge. One end of $\gamma$ is then an end of this chain and carries the measurement of $x$. The other is the other end of the chain, or a neighbor of it, which is a copy in the same context or a vertex of a chain between copies of $x$; its measurement commutes with that of $x$.

\emph{Divisibility.} Consider any quantum $3$-coloring of $F_{22}^{(l)}$. A diamond is the union of two triangles with a common edge, so by~(Q2) its two tips carry the same projectors. Hence all copies of a vertex $x$ of $F_{22}$ carry the same measurement, and assigning it to $x$ gives a quantum $3$-coloring of $F_{22}$ in the same local dimension, which is divisible by four by~(Z6).
\end{proof}

For $l=1,2$ the program \texttt{check\_note.py} finds that the measurements of this coloring commute up to distance $2l+1$ and not beyond.

\subsection{The 165 measurements with fewer edges}\label{sec:165}

We recall the graph $F_{\rm phys}$ and the tower over it from~\cite[Section~6.2 and Appendix~E]{Z}. Let $C=\mathbb Z/3$ be the set of colors, and index the positions $1,2,3$ of a basis by $i\in C$ in this order. Thirteen \emph{difference triples} in $C^3$ are used:
\[
 O=(0,0,0),\qquad\alpha_i=\mathbf1-e_i,\qquad\beta_i=-\alpha_i,\qquad\delta_i=-e_i,\qquad\upsilon_i=(s-i)_{s\in C}\qquad(i\in C).
\]
For $k\in C$ and a difference triple $p$ the word $\omega(k,p)=(k,k+p_0,k+p_1,k+p_2)$ is a word of the basis chart of $S_0(K_3)$ at each basis. Its measurement in basis $b$ is written $p/k$ when the basis is clear. The graph $F_{\rm phys}$ is the physical graph on the measurements of these $39$ words in the six bases. It has $165$ vertices: the three \emph{scalar} measurements, from $p=O$, which occur in all six bases; six measurements on each of the nine planes, from $p=\alpha_i,\beta_i$, each occurring in the two bases of its plane; and $108$ measurements with an outcome of rank one, from $p=\delta_i,\upsilon_i$, each occurring in one basis. In basis $b$ and for each $i\in C$, the nine words $\omega(k,p)$ with $p\in\{O,\alpha_i,\beta_i\}$ take one color on the positions $\{0,i+1\}$ and one color on the other two; for the three values of $i$ these are the end grids of the three wires at $b$.

For $p\ne q$ in $C^3$ we say that $p$ \emph{precedes} $q$ if all coordinates of $q-p$ are $0$ or $1$. Let $\Delta$ be the simplicial complex on the thirteen difference triples whose simplices are the sets $\{p_0,\ldots,p_m\}$ in which $p_a$ precedes $p_{a'}$ whenever $a<a'$; it is a subcomplex of the triangulation of the torus $(\R/3\mathbb Z)^3$ by the ordered simplices of its unit cubes. Any two vertices of a simplex are comparable. It contains the $21$ triangles
\begin{equation}\label{eq:triangles}
\begin{gathered}
 (O,\delta_i,\alpha_i),\quad(\delta_i,\alpha_i,\upsilon_{i+1}),\quad(\alpha_i,\upsilon_{i+1},\upsilon_i),\quad(\alpha_i,\upsilon_i,\beta_i),\\
 (\beta_i,\upsilon_i,\upsilon_{i+2}),\quad(\beta_i,\upsilon_{i+2},\delta_{i+1}),\quad(\beta_i,\delta_{i+1},O)\qquad(i\in C).
\end{gathered}
\end{equation}
Let $\Delta_{21}$ be their union. It has $36$ edges: the nine edges $O\alpha_i$, $\alpha_i\beta_i$, $\beta_iO$, each in one triangle, and $27$ edges that lie in two triangles each. The sum of the $21$ triangles therefore has the three loops $O,\alpha_i,\beta_i,O$ as its boundary modulo two.

The circle space $Z$ of~\cite[Appendix~E]{Z} consists of six copies of $(\R/3\mathbb Z)\times\Delta$, on which the circle $\R/3\mathbb Z$ acts by translation of the first coordinate $t$, and of nine cylinders that join, for each edge $e$ of $K_{3,3}$, the tori over the loops of the two end grids of $e$, with a twist at the complemented edge. The map $Z\to\B(F_{\rm phys})$ sends the point $(t,p)$ of the copy of basis $b$ to the pair of the distributions of the measurement of the word
\[
 \bigl(\lfloor\sigma\rfloor,\lfloor\sigma+p_0\rfloor,\lfloor\sigma+p_1\rfloor,\lfloor\sigma+p_2\rfloor\bigr)\bmod3
\]
in basis $b$ for $\sigma$ uniform in $[t,t+\frac12]$ and for $\sigma$ uniform in $[t+\frac32,t+2]$, and it is constant along the cylinders. The words that occur at $(t,p)$ are words $\omega(k,p')$ with $p'$ a vertex of the smallest simplex of $\Delta$ that contains $p$, and two words in opposite shores differ at all four positions. The tower consists of a point $x_0$ with $p=O$ and a $2$-chain $x_1$ with $\partial x_1=Bx_0$, which exists because the sum of the three loops bounds in $\Delta$ and the cylinder of the complemented edge is attached with a twist.

Let $Z_{21}\subseteq Z$ be the subspace built in the same way from $\Delta_{21}$. It is invariant under the circle action and contains the loops, the cylinders and the $21$ triangles, and $\Delta_{21}$ is connected; the argument of~\cite[Appendix~E]{Z} for the existence of $x_1$ uses only this. So $Z_{21}$ is a circle space with a tower $x_0,x_1$ of length one.

\paragraph{Two subgraphs.}
Let $F_\triangle$ be the spanning subgraph of $F_{\rm phys}$ with the following edges, for each basis: for each difference triple $p$, the three edges among the measurements $p/0,p/1,p/2$; and for each edge $\{p,q\}$ of $\Delta_{21}$, with $p$ preceding $q$, the three edges from $p/k$ to $q/(k+1)$, $k\in C$. These pairs are edges of $F_{\rm phys}$, because the words $\omega(k,p)$ and $\omega(k+1,q)$ differ at all four positions. Let $F_*$ be obtained from $F_\triangle$ by deleting, for each basis and each $i\in C$, the three edges among $\delta_i/0,\delta_i/1,\delta_i/2$.

The following lemma allows the deletion of edges without loss of the tower.

\begin{lemma}[Deleting an edge]\label{lem:link}
Let $uv$ be an edge of a graph $G$, and let $\Lambda$ be the simplicial complex whose vertex set is the disjoint union of $N(v)\setminus\{u\}$, the \emph{left} vertices, and $N(u)\setminus\{v\}$, the \emph{right} vertices, and whose simplices are the sets in which every left vertex is adjacent in $G$ to every right vertex. If $\Lambda$ is contractible, there is an equivariant map $\B(G)\to\B(G-uv)$ that is the identity on $\B(G-uv)$. In particular, every tower over $G$ gives a tower over $G-uv$ on the same circle space.
\end{lemma}
\begin{proof}
A point $(\mu,\nu)$ of $\B(G)$ lies outside $\B(G-uv)$ exactly if $u$ and $v$ lie in opposite shores. Let $\Sigma_{uv}$ be the closure of the set of points with $\mu(u)>0$ and $\nu(v)>0$. Its points are described by the weights $\xi$ of $\mu$ off $u$ and $\eta$ of $\nu$ off $v$: these are nonnegative functions on the left and on the right vertices, respectively, with sums at most one, such that every left vertex in the support of $\xi$ is adjacent to every right vertex in the support of $\eta$. The point lies in $\B(G-uv)$ exactly if one of the two sums equals one. Hence $\Sigma_{uv}$ is the cone, with apex $\xi=\eta=0$, over $\Sigma_{uv}\cap\B(G-uv)$, the set where the larger of the two sums is one; and radial projection maps this set homeomorphically onto the set where the two sums add up to one, which is $\Lambda$. The corresponding set $\Sigma_{vu}$ is the image of $\Sigma_{uv}$ under the involution, and the two meet only in $\B(G-uv)$, because $u$ is not adjacent to itself. So $\B(G)$ arises from $\B(G-uv)$ by attaching two cones over copies of $\Lambda$, which the involution exchanges. If $\Lambda$ is contractible, the base of a cone is a deformation retract of the cone~\cite[Chapter~0]{Hatcher}. Choosing a retraction on $\Sigma_{uv}$ and its conjugate by the involution on $\Sigma_{vu}$ gives the required map. Composing with it turns an equivariant map from a circle space to $\B(G)$ into one to $\B(G-uv)$.
\end{proof}

\begin{theorem}\label{thm:star}
The graphs $F_\triangle$ and $F_*$ are spanning subgraphs of $F_{\rm phys}$ with $732$ and $678$ edges. There is a tower of length one over each of them. Consequently $w^3\ne0$ on their box complexes, $(\cq,\chi)=(3,5)$ for both, and each has a quantum $3$-coloring in local dimension four and none in a smaller local dimension.
\end{theorem}
\begin{proof}
\emph{Numbers of edges.} Fix a basis $b$. The definition of $F_\triangle$ lists $13\cdot3+36\cdot3=147$ pairs of words of $b$. They are distinct pairs, and distinct words of one basis have distinct measurements, so they give a set $E_b$ of $147$ edges. No edge of $\Delta_{21}$ joins two of the triples $\alpha_i,\beta_i$ with different indices, as the list~\eqref{eq:triangles} shows. Hence the edges of $E_b$ both of whose ends occur in further bases are the edges inside the three end grids. In the end grid with index $i$ the definition gives the three edges among the scalar measurements, the six edges among the $\alpha_i/k$ and among the $\beta_i/k$, and nine edges from the three edges $O\alpha_i$, $\alpha_i\beta_i$, $\beta_iO$ of $\Delta_{21}$. These are $18$ edges, and the grid $K_3^{\times2}$ has $18$ edges, so every end grid is complete in $E_b$. There are $3+3\cdot15=48$ such edges in $E_b$, and the other $99$ edges of $E_b$ have an end with an outcome of rank one, which belongs to basis $b$ only. Two bases that share a plane give the same $18$ edges on its grid. Altogether $F_\triangle$ has $3+9\cdot15+6\cdot99=732$ edges. The deleted edges join measurements that belong to one basis, so they are $6\cdot3\cdot3=54$ distinct edges, and $F_*$ has $678$ edges. The same count shows that the subgraph of $F_\triangle$ induced on the $39$ measurements of basis $b$ has edge set $E_b$.

\emph{A tower over $F_\triangle$.} We show that the map $Z\to\B(F_{\rm phys})$ sends $Z_{21}$ into $\B(F_\triangle)$. Let $(t,p)$ be a point of a copy of $(\R/3\mathbb Z)\times\Delta_{21}$, and let $\omega(k,p')$ and $\omega(k',q')$ be words in opposite shores of its image. Then $p'$ and $q'$ are vertices of a triangle of $\Delta_{21}$ and the two words differ at all four positions. If $p'=q'$, then $k\ne k'$ and the edge is one of the first kind. Otherwise we may assume that $p'$ precedes $q'$. Position $0$ gives $k'\ne k$. Some coordinate of $q'-p'$ is $1$, and the corresponding position gives $k'\ne k-1$. So $k'=k+1$ and the edge is of the second kind. The points of the cylinders have the same images as points of the copies. With the tower $x_0,x_1$ on $Z_{21}$ this is a tower of length one over $F_\triangle$.

\emph{A tower over $F_*$.} We delete the $54$ edges one by one and apply Lemma~\ref{lem:link} each time. Consider first the edge between $u=\delta_0/0$ and $v=\delta_0/1$ of a basis $b$ in $F_\triangle$. The edges of $\Delta_{21}$ at $\delta_0$ are $\delta_0O$, $\delta_0\alpha_0$, $\delta_0\upsilon_1$ and $\beta_2\delta_0$, each written with the preceding triple first. So the neighbors of $\delta_0/k$ are $\delta_0/(k\pm1)$, $O/(k+1)$, $\alpha_0/(k+1)$, $\upsilon_1/(k+1)$ and $\beta_2/(k-1)$, all in basis $b$. The left and right vertices of $\Lambda$ and their adjacencies in $F_\triangle$ are given by the following table, with the left vertices as rows and an entry $1$ for adjacent pairs.
\[
\begin{array}{c|ccccc}
  &\alpha_0/1&\beta_2/2&\delta_0/2&O/1&\upsilon_1/1\\ \hline
\alpha_0/2&1&0&0&1&1\\
\beta_2/0&0&1&0&1&1\\
\delta_0/2&0&0&0&0&0\\
O/2&0&0&0&1&0\\
\upsilon_1/2&0&0&0&0&1
\end{array}
\]
The entries follow from the definition of $E_b$: among the triples $O,\alpha_0,\beta_2,\delta_0,\upsilon_1$ the edges of $\Delta_{21}$ are, besides those at $\delta_0$, the edges $O\alpha_0$, $\beta_2O$, $\upsilon_1\alpha_0$ and $\beta_2\upsilon_1$, again with the preceding triple first.

If the right neighbors of a left vertex $a$ of $\Lambda$ are among those of another left vertex $a'$, then every simplex that contains $a$ remains a simplex when $a'$ is added, so the link of $a$ is a cone and removing $a$ does not change the homotopy type; the same holds with left and right exchanged. In the table, the last three rows are dominated in this sense by the first row. After their removal the first three columns are dominated by the fourth. The remaining two rows and two columns have all entries equal to $1$ and span a simplex. So $\Lambda$ is contractible.

The set $E_b$ is invariant under adding a constant to all colors and under cyclic permutations of the positions $1,2,3$, which permute the difference triples and the triangles~\eqref{eq:triangles} and preserve the relation of preceding. Hence for each of the $54$ edges the complex $\Lambda$ formed in $F_\triangle$ is given by the same table up to these symmetries and up to an exchange of left and right. When some of the $54$ edges have already been deleted, the only possible change of the complex $\Lambda$ of a further one is that the vertex in the third row, in the third column, or both, are absent, because the other vertices are not ends of deleted edges and all entries of that row and column are zero. The reduction above still applies. By Lemma~\ref{lem:link}, applied $54$ times, the tower over $F_\triangle$ gives a tower over $F_*$.

\emph{Chromatic numbers and local dimension.} By~(T1), $w^3\ne0$ on both box complexes and both chromatic numbers are at least five. They are at most five because $\chi(F_{\rm phys})=5$ by~(Z7). The measurements themselves form a quantum $3$-coloring in local dimension four, and the three scalar measurements form a triangle in $F_*$, so $\cq=3$ by~(Q3). Since $w^2\ne0$, every quantum $3$-coloring has even local dimension by~(Z5), and local dimension two is excluded by~(Z4).
\end{proof}

\subsection{Chromatic number five at a prescribed compatibility radius}\label{sec:FR}

Let $R\ge2$ and $\lambda=R-2$. The graph $F_R$ consists of six \emph{charts} and nine wires. The chart of basis $b$ has the $39$ words $\omega(k,p)$ of basis $b$ as vertices and the $147$ pairs of words that define $E_b$ as edges; the six charts are disjoint. For an edge $e=ik$ of $K_{3,3}$, layer $0$ of its wire is the end grid of $e$ in chart $i$ and layer $\lambda+1$ is the end grid of $e$ in chart $k$, where the cell $(x,y)$ is the word that takes $x$ on $A_{b,e}$ and $y$ on $A_{b,e}^c$. An interior layer $\ell\in\{1,\ldots,\lambda\}$ consists of the six vertices $(e,\ell,x,y)$ with $y\in C$ and $x\in C\setminus\{m_\ell\}$, where $m_\ell=0$ for odd $\ell$ and $m_\ell=1$ for even $\ell$. Two cells $(x,y)$ and $(x',y')$ in equal or consecutive layers of a wire are adjacent when $x\ne x'$ and $y\ne y'$; inside an end grid these are the edges of the chart. A word carries its measurement, and $(e,\ell,x,y)$ carries the measurement that returns $x$ on the plane of $e$ and $y$ on its complement.

\begin{theorem}\label{thm:FR}
Let $R\ge2$. The graph $F_R$ has $234+54(R-2)$ vertices, which carry the $165$ measurements of $F_{\rm phys}$. These measurements form a quantum $3$-coloring of $F_R$ in local dimension four whose compatibility radius is exactly $R$. There is a tower of length one over $F_R$, so $w^3\ne0$ on $\B(F_R)$ and $(\cq,\chi)=(3,5)$. No quantum $3$-coloring of $F_R$ has local dimension less than four.
\end{theorem}
\begin{proof}
The order is $6\cdot39+9\cdot6\lambda$.

\emph{Coloring and upper bound.} Sending a vertex to its measurement is a homomorphism $F_R\to F_\triangle$: the edges of a chart go to edges of $E_b$, and an edge of a wire goes to an edge of the grid of its plane, which is complete in $F_\triangle$. So the measurements form a quantum $3$-coloring and $\chi(F_R)\le5$. A chart contains the triangle of the three constant words, so $\cq(F_R)=3$.

\emph{Compatibility.} Every vertex of the chart of $b$ and every vertex of a wire at $b$, including the end grid at its other end, carries a measurement that is diagonal in basis $b$. A path between two charts runs through a wire and has at least $\lambda+1$ edges. As in the proof of~\cite[Lemma~2.2]{Z}, a path with at most $\lambda+2$ edges therefore has all its vertices in the chart of one basis $b$ or on wires at $b$, and the measurements at its ends commute.

\emph{The radius is not larger.} Take the wire of $e=24$, for which $A_{2,e}=A_2$ and $A_{4,e}=A_1$. Choose colors $z_0,\ldots,z_{\lambda+1}$ with $z_\ell\ne z_{\ell+1}$ for $0\le\ell\le\lambda$ and $z_\ell\ne m_\ell$ for $1\le\ell\le\lambda$; this is possible because every interior layer leaves two colors. The cells $(z_\ell,z_\ell)$ of the layers $\ell=0,\ldots,\lambda+1$ form a path with $\lambda+1$ edges from the constant word $z_0$ of chart $2$ to the constant word $z_{\lambda+1}$ of chart $4$. In chart $2$ the word $\omega(z_0-1,\delta_1)$ is adjacent to the constant word $z_0$, since $\delta_1$ precedes $O$, and one outcome of its measurement is the projection $Q^2_2$ onto the ray of $(0,0,1,0)$. In chart $4$ the word $\omega(z_{\lambda+1}-1,\delta_0)$ is adjacent to the constant word $z_{\lambda+1}$, and one outcome of its measurement is the projection $Q^4_1$ onto the ray of $(1,0,-1,0)$. These two projections do not commute, and the two words are joined by a path with $\lambda+3=R+1$ edges.

\emph{A tower.} We define an equivariant map $Z_{21}\to\B(F_R)$. On the copy of basis $b$ it is the map of Section~\ref{sec:165}, with the words regarded as vertices of the chart of $b$; opposite shores are completely adjacent in the chart by the argument for $F_\triangle$. It remains to extend it over the cylinders. Write $Q=K_3^{\times2}$. Two homomorphisms $f,g$ from $Q$ to a graph $G$ are \emph{compatible} if $f(q)\sim g(q')$ whenever $q\sim q'$ in $Q$. For compatible $f,g$ the maps
\[
 (\mu,\nu)\mapsto\bigl((1-a)f_*\mu+a\,g_*\mu,\ (1-a)f_*\nu+a\,g_*\nu\bigr),\qquad0\le a\le1,
\]
form an equivariant homotopy from $\B(f)$ to $\B(g)$, where $f_*$ denotes the image of a measure under $f$. Fix a wire. Two consecutive layers contain every cell of $Q$ at least once, because the end grids are complete and consecutive interior layers omit different rows. A \emph{selection} for two consecutive layers assigns to every cell of $Q$ a vertex with the same coordinates in one of the two layers. Every selection is a homomorphism $Q\to F_R$, and two selections for the same two layers are compatible. For an interior layer $\ell$ let $c^-_\ell$ and $c^+_\ell$ be the selections for the layers $\ell-1,\ell$ and for the layers $\ell,\ell+1$ that use layer $\ell$ for its six cells. They are compatible: if $q\sim q'$ and neither $c^-_\ell(q)$ nor $c^+_\ell(q')$ lies in layer $\ell$, then $q$ and $q'$ both lie in the row that layer $\ell$ omits, which is impossible for adjacent cells; otherwise the two vertices lie in equal or consecutive layers and are adjacent. Let $c_0$ and $c_{\lambda+1}$ be the selections given by the two end grids. In the sequence $c_0,c^-_1,c^+_1,c^-_2,\ldots,c^+_\lambda,c_{\lambda+1}$ consecutive members are compatible, so the homotopies above combine to an equivariant homotopy from $\B(c_0)$ to $\B(c_{\lambda+1})$.

On the torus over the loop of the end grid of $e$ in the copy of basis $i$, the map of Section~\ref{sec:165} is $\B(c_0)$ composed with a map to $\B(Q)$, namely the pair of the distributions of the cell $(\lfloor\sigma\rfloor,\lfloor\sigma+\vartheta\rfloor)$, where $\vartheta$ is the coordinate along the loop. We extend it over the cylinder by the homotopy just constructed. At the other end the result is $\B(c_{\lambda+1})$ composed with the same map to $\B(Q)$, which is the map of Section~\ref{sec:165} on the torus in the copy of basis $k$: the attaching maps of the cylinders of $Z$ are chosen in~\cite[Appendix~E]{Z} so that the words at the two ends have the same color on the plane of $e$ and the same color on its complement, also at the complemented edge. This defines an equivariant map $Z_{21}\to\B(F_R)$, and with $x_0,x_1$ a tower of length one over $F_R$. Hence $w^3\ne0$ on $\B(F_R)$ and $\chi(F_R)\ge5$ by~(T1). The assertion on the local dimension follows from~(Z5) and~(Z4) as in Theorem~\ref{thm:star}.
\end{proof}

\subsection{A physical family that starts with the 165 measurements}\label{sec:Hj}

By~\cite[Corollary~6.3]{Z} the tagged graphs $S^{j-1}_\lambda(F_{\rm phys})$ have chromatic number at least $2j+3$; their order is doubly exponential in $j$. We describe a sparse family of physical graphs with the same first level.

\paragraph{Presentations and weights.}
A measurement \emph{at level $g$} is a projective measurement with three outcomes on $(\C^4)^{\otimes g}$. For a color $c$, the \emph{padded scalar} $c$ at level $g$ is the measurement that returns $c$ with certainty; it has \emph{weight} $0$. At level one we use the $165$ measurements of $F_{\rm phys}$: the three scalar ones are the padded scalars, and the other $162$ have weight $1$. A \emph{template node} at level $g\ge2$ consists of a template at a basis $i$, with parts $D_1,\ldots,D_p$, and of a \emph{child} $N^r$ for each part, which is a padded scalar at level $g-1$, a template node at level $g-1$ or, if $g=2$, one of the $162$ other measurements. It presents the measurement
\begin{equation}\label{eq:node}
 M_a=\sum_{r=1}^p\Bigl(\sum_{s\in D_r}Q^i_s\Bigr)\otimes N^r_a .
\end{equation}
A child is \emph{internal} if it is not a padded scalar. The weight of a template node is the sum of the weights of its children, plus one if at most one child is internal. A \emph{presentation} at level $g$ is a padded scalar or a template node, or, for $g=1$, one of the $165$ measurements.

For $j\ge1$ let $H_j$ be the physical graph on all measurements at level $j$ that have a presentation of weight at most $j$. Thus $H_1=F_{\rm phys}$.

\begin{lemma}\label{lem:bottom}
Let $\mathcal B$ be a connected graph in which every vertex lies in a triangle, and let $(P_{(y,v),a})$ be a quantum $3$-coloring on $\C^d$ of the categorical product $K_3\times\mathcal B$, in which $(y,v)\sim(y',v')$ when $y\ne y'$ and $v\sim v'$. Then there is a projection $L$ that commutes with all projectors of the coloring such that on the range of $L$ the measurement at $(y,v)$ does not depend on $y$, and such that for every triangle $\mathcal T$ of $\mathcal B$ the projection $L$ is the coordinate projection of the second coordinate for the restriction of the coloring to $K_3\times\mathcal T=K_3^{\times2}$. If $L\ne0$, the measurements on its range form a quantum $3$-coloring of $\mathcal B$ in local dimension $\rank L$.
\end{lemma}
\begin{proof}
Let $v$ be a vertex and $\mathcal T$ a triangle that contains it. By~(Z2), applied to $K_3\times\mathcal T$, there is a decomposition $\C^d=V_1\oplus V_2$ that reduces the projectors at the vertices of $K_3\times\mathcal T$, such that on $V_2$ the measurement at $(y,v')$, $v'\in\mathcal T$, depends only on $v'$, and on $V_1$ only on $y$, with the edge rule of $K_3$. The operator $L_v=\sum_aP_{(0,v),a}P_{(1,v),a}$ is therefore the projection onto $V_2$. As it is defined without reference to $\mathcal T$, the space $V_2$ is the same for all triangles at $v$, and $L_v=L_{v'}$ for $v'\in\mathcal T$. Moreover $L_v$ commutes with the projectors at $(y,v)$, and on the range of $I-L_v$ the three projections $P_{(z,v),a}$, $z\in K_3$, of one color sum to the identity by~(Q2).

Let $v\sim w$ and let $\xi$ be a vector in the range of $L_v$ on which the common measurement at $(y,v)$, $y\in K_3$, returns the color $a$. For every $z$ choose $y\ne z$; then $(y,v)\sim(z,w)$ and $P_{(z,w),a}\xi=0$. So $\sum_zP_{(z,w),a}$ annihilates $\xi$. This sum commutes with $L_w$ and is the identity on the range of $I-L_w$, hence $(I-L_w)\xi=0$. Such vectors $\xi$ span the range of $L_v$, so $L_v\le L_w$, and $L_v=L_w$ by symmetry. As $\mathcal B$ is connected, all $L_v$ equal one projection $L$, which commutes with all projectors. On its range the measurement at $(y,v)$ does not depend on $y$, and for $v\sim w$ the edge rule for $(0,v)$ and $(1,w)$ gives the edge rule for $v$ and $w$.
\end{proof}

\begin{theorem}\label{thm:Hj}
Let $j\ge1$.
\begin{enumerate}
\item[(a)] The measurements that are the vertices of $H_j$ form a quantum $3$-coloring of $H_j$ in local dimension $4^j$, and $\cq(H_j)=3$.
\item[(b)] There is a tower of length $j$ over $H_j$. Hence $w^{2j+1}\ne0$ on $\B(H_j)$ and $\chi(H_j)\ge2j+3$.
\item[(c)] There is an absolute constant $c_0$ with $|V(H_j)|\le c_0^{\,j}$ and $|E(H_j)|\le c_0^{\,2j}$.
\item[(d)] Every quantum $3$-coloring of $H_j$ has local dimension divisible by $2\cdot4^{j-1}$ and at least $4^j$.
\item[(e)] For $j\ge2$, $\chi(H_j)\le2\chi(H_{j-1})-1$. The graph $H_2$ has $4{,}330{,}533$ vertices and $7\le\chi(H_2)\le9$.
\end{enumerate}
\end{theorem}
\begin{proof}
\emph{Part~\textup{(a)}.} Adjacent vertices are orthogonal at every outcome. The three padded scalars have weight $0$ and form a triangle, so~(Q3) applies.

\emph{Part~\textup{(c)}.} In a presentation, call the template nodes and the nonscalar measurements at level one \emph{internal nodes}; with the relation of being a child they form a tree. An internal node without internal children contributes one to the weight, as does an internal node with exactly one internal child, and no other node contributes. If the weight is at most $j$, the tree has at most $j$ leaves, hence at most $j-1$ nodes with two or more children, and at most $2j-1$ nodes in all. Each node is one of $162$ measurements or one of nineteen templates with at most three children, each a padded scalar or an internal node. So there are at most $c_0^{\,j}$ presentations of weight at most $j$, and the number of edges is at most the square of the number of vertices.

\emph{Part~\textup{(b)}: a tower on a space of presentations.} Let $Y_1=Z_{21}$ and $Y_g=S(Y_{g-1})$ for $g\ge2$. Let $O_1\subset Y_1$ be the orbit of the point $x_0$, which we take in the copy of basis $6$, and let $O_g\subset Y_g$ consist of the constant tuples in the block of basis $6$ on the points of $O_{g-1}$. We give weights to continuous presentations of points of $Y_g$. At level one a presentation is the point itself, of weight $0$ and not internal if it lies on $O_1$, and of weight $1$ and internal otherwise. At level $g\ge2$ a presentation is a point of $O_g$, of weight $0$ and not internal; or a basis node, with a template and children at level $g-1$, whose weight is the sum of the weights of the children plus one if at most one child is internal; or a wire node, with a first child at level $g-1$ and a background in $O_{g-1}$, whose weight is the weight of the first child plus one. Basis nodes and wire nodes are internal. Let $Y_{g,m}\subseteq Y_g$ be the set of points with a presentation of weight at most $m$. These sets are invariant under the circle action. They are compact polyhedra: regarding a point of $O_1$ also as an internal presentation of weight $1$ does not change the sets $Y_{g,m}$, because it raises the sum of the weights of the children of a node by one and lowers the added unit by at most one; and with this convention $Y_{g,m}$ is a finite union of images of products of copies of $Y_1$, of $O_1$ and of intervals under piecewise linear maps.

We claim that $Y_g$ carries a tower $x_0,\ldots,x_g$ with $x_0$ a point of $O_g$ and $x_m\in C_{2m}(Y_{g,m};\Ftwo)$ for $1\le m\le g$. For $g=1$ this is the tower on $Z_{21}$, since every point of $Y_1$ has weight at most one. For the step from $g$ to $g+1$ we take the tower of~(T2) and examine the three kinds of terms of $y_m$ in the proof of Lemma~\ref{lem:supports}. In a term of correction order zero with at least two factors of positive indices $i_\nu$, which sum to $m$, the children have weights at most $i_\nu$, and the other children are points of $O_g$. If at least two children are internal, the weight is at most $m$. If at most one is internal, one of the factors of positive index contributes a point of $O_g$, of weight $0\le i_\nu-1$, and the weight is at most $(m-1)+1$. In a term of correction order $r\ge1$ the weights of the children sum to at most $m-r$, and the weight is at most $m-r+1\le m$. The points of $BKx_{m-1}$ and $KBx_{m-1}$ are presented by wire nodes over points of $Y_{g,m-1}$, with weight at most $m$. This proves the claim. For $g=j$ it gives a tower of length $j$ on the circle space $Y_{j,j}$.

\emph{Part~\textup{(b)}: the map to $\B(H_j)$.} Let $\Phi_1:Y_1\to\B(F_{\rm phys})$ be the map of Section~\ref{sec:165} and let $\Phi_g:Y_g\to\B(S^{g-1}_0(F_{\rm phys}))$ arise from it by $g-1$ applications of~(T3) with $\lambda=0$. The $(g-1)$-fold iterate of~\eqref{eq:strategy}, starting from the measurements of $F_{\rm phys}$, assigns to every vertex of $S^{g-1}_0(F_{\rm phys})$ a measurement at level $g$, and adjacent vertices receive measurements that are orthogonal at every outcome. We show, by induction on the presentation, that if a point $z\in Y_g$ has a presentation $\pi$ of weight $m$, then the measurement of every vertex in a shore of $\Phi_g(z)$ has a presentation $\pi'$ of weight at most $m$; moreover $\pi'$ is internal if $\pi$ is, except that $\pi'$ may be a padded scalar when $\pi$ is a point of $Y_1$ off $O_1$, in which case its weight is $0=m-1$.

For points of $Y_1$ this holds because $\Phi_1$ maps $O_1$ to pairs of distributions on the scalar measurements. For a point of $O_g$, $g\ge2$, the vertices in the shores are constant words in basis $6$ on vertices in the shores one level below, by the property in~(T3), and their measurements are padded scalars. For a basis node with children $\pi_1,\ldots,\pi_p$, a vertex in a shore is a word of the chart of its basis whose entries are constant on the parts and come from the shores of the children, so by~\eqref{eq:strategy} its measurement is presented by the template node with children $\pi'_1,\ldots,\pi'_p$. If every internal child $\pi_r$ has $\pi'_r$ internal, the number of internal children is unchanged and the weight does not increase. Otherwise some weight drops by one, the added unit is at most one, and again the weight does not increase. For a wire node at the edge $e=ik$ with first child $\pi_1$, a vertex in a shore is a pair in layer $0$ or $1$ of the wire of $e$ in $S_0$, that is, a word of an end grid of $e$; by~\eqref{eq:strategy} and~\eqref{eq:planes} its measurement is $P_e\otimes N+(I-P_e)\otimes c$ with $N$ presented by $\pi'_1$ and $c$ a padded scalar. It is presented by the ternary template at basis $i$ with the part $A_{i,e}$, with child $\pi'_1$ on this part and the padded scalar $c$ on the two singletons, of weight at most the weight of $\pi_1$ plus one.

Consequently, for $z\in Y_{j,j}$ the measurements of the vertices in the two shores of $\Phi_j(z)$ are vertices of $H_j$, and those of opposite shores are adjacent in $H_j$. Taking the images of the two distributions under the map from vertices to measurements gives an equivariant map $Y_{j,j}\to\B(H_j)$. With the tower on $Y_{j,j}$ this is a tower of length $j$ over $H_j$, and~(T1) gives the rest of~(b).

\emph{Part~\textup{(d)}.} Fix a quantum $3$-coloring $(P_{M,a})$ of $H_j$ on $\C^d$. For a color $y$, a basis $i$, a position $s\in\{1,2,3\}$ and a measurement $N$ at level $g-1$, let $[i,s;y,N]$ be the measurement $Q^i_s\otimes N+(I-Q^i_s)\otimes y$ at level $g$. It is presented by a template node with a ternary template at basis $i$ that has the part $\{s\}$, with $N$ on this part and the padded scalar $y$ on the other two. For a sequence $\tau=((i_1,s_1),\ldots,(i_h,s_h))$ with $0\le h\le j-1$ put
\[
 [\tau;y,N]=[i_1,s_1;y,[i_2,s_2;y,\ldots[i_h,s_h;y,N]\ldots]]
\]
for a measurement $N$ at level $j-h$. Each node has at most one internal child, so $[\tau;y,N]$ has a presentation of weight $h$ plus the weight of $N$. In particular $F_\tau(y,x)=[\tau;y,x]$, with the padded scalar $x$, has weight $h$, and $F_\tau$ is a homomorphism $K_3^{\times2}\to H_j$. Let $\Pi_\tau$ be the coordinate projection of the coordinate $x$ for the coloring pulled back along $F_\tau$; see~(Z2). For the empty sequence $\Pi_\tau=I$.

Let $h\le j-2$. We claim that $\rank \Pi_{\tau'}=\frac14\rank \Pi_\tau$ for every sequence $\tau'=(\tau,(i,s))$. This is the step of the proof of~\cite[Theorem~4.1]{Z}, which we summarize. Inserting a ternary template node at basis $i$ with three scalar children $x_1,x_2,x_3$ as the innermost measurement gives a homomorphism $K_3^{\times4}\to H_j$ in the variables $y,x_1,x_2,x_3$, with presentations of weight $h+1$. For $x_1=x_2=x_3=x$ the measurement is $F_\tau(y,x)$, so the coordinate projections of the three parts sum to $\Pi_\tau$. The projection of a part $\{s\}$ is $\Pi_{(\tau,(i,s))}$, whichever of the two templates with this part is used, because the measurement with $x$ at position $s$ and $y$ elsewhere is the same. So the projections $T^i_s=\Pi_{(\tau,(i,s))}$, $s=1,2,3$, are pairwise orthogonal, the projection of the part $A_r$ is $D^i_r=\Pi_\tau-T^i_s-T^i_t$ for $\{r,s,t\}=\{1,2,3\}$, and with $T^i_0=\Pi_\tau-\sum_sT^i_s$ one has $D^i_r=T^i_0+T^i_r$. For an edge $e=ik$ of $K_{3,3}$ the measurements with $x$ on $A_{i,e}$ in basis $i$ and with $x$ on $A_{k,e}$ in basis $k$, and $y$ on the other positions, are equal by~\eqref{eq:planes}; they are one vertex of $H_j$. Hence $D^i_r=D^k_{r'}$ if $A_{i,e}=A_r$ and $A_{k,e}=A_{r'}$, and $D^3_3=T^6_1+T^6_2=\Pi_\tau-D^6_3$ for $e=36$. On the range of $\Pi_\tau$ the involutions $2D^i_r-\Pi_\tau$ thus satisfy the relations that~\cite{Z} derives for the tagged graphs by transport along the wires, and the argument given there with~(Z3) shows that $T^i_0,\ldots,T^i_3$ have equal rank. This proves the claim. The restriction $h\le j-2$ ensures that the inserted node lies at level at least two, where its children may be arbitrary scalars.

Applying the claim $j-1$ times gives a sequence $\tau$ of length $j-1$ with $\rank \Pi_\tau=d/4^{j-1}$; in particular $4^{j-1}$ divides $d$. For this $\tau$ and a vertex $v$ of $F_{\rm phys}$, the measurement $[\tau;y,v]$ has weight at most $j$, and $(y,v)\mapsto[\tau;y,v]$ is a homomorphism $K_3\times F_{\rm phys}\to H_j$. The graph $F_{\rm phys}$ is connected and each of its vertices lies in a triangle: the three measurements $p/0,p/1,p/2$ form a triangle; a measurement with at most two colors is adjacent to the scalar of the third color; and if the word $\omega(k,\upsilon_i)$ has the color $c\ne k$ at the position $s\ge1$, it differs at all four positions from the word with $c-1$ at $s$ and $c$ elsewhere, which is $\omega(c,\delta_{s-1})$ and has two colors. Apply Lemma~\ref{lem:bottom} to the pulled-back coloring of $K_3\times F_{\rm phys}$. For the triangle of the scalar measurements the homomorphism is $F_\tau$, so $L=\Pi_\tau$, of rank $d/4^{j-1}>0$. Hence $F_{\rm phys}$ has a quantum $3$-coloring in local dimension $d/4^{j-1}$. By~(Z7), (Z5) and~(Z4) this number is even and at least four, as in Theorem~\ref{thm:star}.

\emph{Part~\textup{(e)}.} Let $j\ge2$. A child of a template node of weight at most $j$ has weight at most $j-1$. Indeed, internal children have weight at least one; so if a child had weight at least $j$, then either a second child would be internal or the unit for a node with at most one internal child would be added, and in both cases the node would have weight at least $j+1$. So every vertex $M$ of $H_j$ is a padded scalar or has the form~\eqref{eq:node} with children that present measurements $N^r\in V(H_{j-1})$. For $N\in V(H_{j-1})$ let $D_N(M)$ be the sum of the projections $\sum_{s\in D_r}Q^i_s$ over the parts $D_r$ whose child presents the measurement $N$; for the padded scalar $c$ at level $j$ let $D_N(M)=I$ if $N$ is the padded scalar $c$ at level $j-1$ and $D_N(M)=0$ otherwise. Then $M_a=\sum_ND_N(M)\otimes N_a$. The family $(D_N(M))_N$ depends only on $M$: a real linear functional on triples of matrices that separates the finitely many elements of $V(H_{j-1})$, applied to the second tensor factor of $M$, gives a self-adjoint operator whose spectral projections are the $D_N(M)$. If $M\sim M'$ in $H_j$, then
\[
 0=\sum_a\tr(M_aM'_a)=\sum_{N,N'}\tr\bigl(D_N(M)D_{N'}(M')\bigr)\sum_a\tr(N_aN'_a),
\]
and all terms are nonnegative. So $D_N(M)D_{N'}(M')\ne0$ implies $N_aN'_a=0$ for all $a$, that is, $N\sim N'$ in $H_{j-1}$. Let $\kappa$ be a proper coloring of $H_{j-1}$ with $n=\chi(H_{j-1})$ colors, and let $\kappa(M)$ be the measurement with labels in $[n]$ that assigns the label $l$ to the sum of the $D_N(M)$ with $\kappa(N)=l$. It is a vertex of the graph $\mathcal P_n$ of~(Z1), and by the preceding observation $M\sim M'$ implies that $\kappa(M)$ and $\kappa(M')$ are orthogonal at every label. So $\kappa$ induces a homomorphism $H_j\to\mathcal P_n$, and $\chi(H_j)\le2n-1$ by~(Z1). For $j=2$ this gives $\chi(H_2)\le9$, as $\chi(H_1)=5$, and $\chi(H_2)\ge7$ by~(b).

\emph{The order of $H_2$.} Put $n=165$ and $b=162$, the numbers of all and of the nonscalar measurements at level one, and write $(n)_k=n(n-1)\cdots(n-k+1)$. A template node at level two with $k$ nonscalar children has weight $k+1$ if $k\le1$ and weight $k$ otherwise, so its weight is at most two exactly if at most two of its children are nonscalar. By the preceding paragraph, a vertex $M$ of $H_2$ is determined by the projections $D_N(M)$ of the measurements $N$ at level one that occur in it, and we count the vertices by the number of these measurements. If there is one, then $M=I\otimes N$; this is a padded scalar, or it is presented by the binary template with both children equal to a nonscalar $N$. This gives $n$ vertices. If there are two, $N$ and $N'$, then $D_N$ and $D_{N'}$ are complementary sums of parts of a template. The pairs of complementary projections that arise are the $18$ pairs $\{Q^i_s,I-Q^i_s\}$ with $s\ge1$ and the nine pairs $\{P_e,I-P_e\}$; they are $27$ distinct pairs by Section~\ref{sec:square}. In a ternary template one of the two measurements occupies two parts, so the weight is at most two exactly if $N$ or $N'$ is scalar; the binary template, which gives the pair of the plane of $e=36$, admits any two measurements. This gives $27\bigl((n)_2-(b)_2\bigr)+(b)_2$ vertices. If there are three, they occupy the three parts of one of the eighteen ternary templates, whose partitions are distinct, and at least one of them is scalar; this gives $18\bigl((n)_3-(b)_3\bigr)$ vertices. More than three cannot occur. In total
\[
 n+27\bigl((n)_2-(b)_2\bigr)+(b)_2+18\bigl((n)_3-(b)_3\bigr)=165+27\cdot978+26{,}082+18\cdot237{,}660=4{,}330{,}533 .\qedhere
\]
\end{proof}

The local dimensions of the quantum $3$-colorings of the tagged graphs $G_{j,R}$ are exactly the multiples of $4^j$~\cite[Theorem~4.1]{Z}. For $H_j$ with $j\ge2$, part~(d) leaves open whether local dimensions such as $6\cdot4^{j-1}$ occur.
